\chapter[Revisão bibliográfica]{Revisão bibliográfica e delimitação da contribuição}
\label{chap:revisao}

\section{Estratégia de revisão e acervo consultado}

Esta revisão organiza os antecedentes necessários para transformar o
problema apresentado na introdução em uma investigação verificável.
O corte bibliográfico é 11 de setembro de 2026. A busca foi dirigida por
quatro eixos: polarizações de campos massivos, resposta de PTA,
inferência de massa e efeitos de compressão e verossimilhança. Foram
combinadas buscas por títulos e identificadores com expressões relativas
a frequência, covariância, priori e gravidade massiva. As referências
dos trabalhos centrais, suas páginas de código e registros editoriais
foram examinados para identificar antecedentes que uma busca somente
por palavras-chave poderia omitir.

O acervo inicial de C03 reuniu 22 PDFs com versões explícitas do arXiv, mantidos no
acervo local \path{literature/papers/}. O catálogo
\path{literature/catalog.json} registra URLs versionadas, datas,
identificadores, tamanho e SHA-256; o protocolo e a matriz de
originalidade acompanham o repositório em \path{docs/literatura/}.
Fixar a versão permite conferir a fonte realmente utilizada, mesmo
quando o artigo recebe uma revisão posterior. Os arquivos foram
verificados quanto ao formato PDF, número de páginas e integridade.
Essas cópias destinam-se
à leitura local e não são redistribuídas automaticamente no Git; as
licenças e a disponibilidade das versões editoriais são questões
distintas do acesso ao preprint. Em C06 foram acrescentadas quatro
referências metodológicas de calibração e evidências bayesianas,
totalizando 26 PDFs e 600 páginas. A leitura dirigida desses métodos
acompanha sua aplicação na etapa de inferência. Os complementos de C07
elevam o acervo a 37 PDFs e 839 páginas. A auditoria dirigida de
11 de setembro acrescenta seis antecedentes à matriz, que passa a
23 entradas; seus registros de busca e localizadores de leitura são
preservados separadamente da busca inicial de C03.

A consulta não constitui uma revisão sistemática exaustiva. O catálogo
distingue leitura dirigida de métodos, consulta anterior e triagem de
resumos. O download não implica leitura integral, assim como a leitura
de um código não implica reprodução de suas saídas. Artigos publicados,
preprints e apresentações são identificados separadamente. Um resultado
negativo da busca significa que determinada análise não foi localizada
nas fontes examinadas, e não que ela inexista. Essa delimitação é
especialmente relevante para trabalhos recentes e será retomada antes
da submissão do manuscrito.

\section{Do TG às teorias de gráviton massivo}

\subsection{Resultado histórico e objeto de reprodução}

O TG de Wayne Leonardo Silva de Paula investiga polarizações em uma
teoria de gravitação com termo de massa, utilizando a linearização em
torno de uma geometria de referência \autocite{tg2003}. Seus capítulos
4--6 articulam equações de propagação, componentes de curvatura e
comparação entre frequências; o apêndice computacional documenta parte
do cálculo. O artigo de de Paula, Miranda e Marinho publicou o resultado
central em 2004 \autocite{depaula2004}. Assim, a obtenção das formas de
polarização e a observação de maior relevância cinemática em baixas
frequências são antecedentes, não hipóteses inéditas desta dissertação.

A reprodução terá função de estabelecer convenções e identificar o
alcance das aproximações. Devem ser explicitadas assinatura da métrica,
convenção de Riemann, unidades e distinção entre a perturbação métrica e
seu traço invertido. A comparação entre amplitudes de curvatura também
depende da normalização escolhida para as amplitudes métricas. Supor
componentes métricas comparáveis é uma condição de comparação; não
determina a distribuição de amplitudes emitida por fontes astrofísicas.
Portanto, a atualização numérica das tabelas históricas é um benchmark
cinemático, e não uma previsão autossuficiente de detectabilidade.

\subsection{Vínculos, estabilidade e limite de validade}

Visser introduz uma formulação com métrica de fundo que permite um
termo massivo e modifica as equações de campo
\autocite{visser1998}. Sua relação com o termo de Fierz--Pauli precisa
ser examinada no nível dos vínculos, não apenas pela presença de uma
relação de dispersão semelhante. A literatura sobre teorias não
Fierz--Pauli evidencia que um modo escalar adicional pode acompanhar
problemas de sinal cinético; uma contagem de soluções não demonstra
estabilidade \autocite{park2010}. O modelo histórico será, por isso,
utilizado como referência de reprodução com seu domínio declarado.

No campo livre de Fierz--Pauli em quatro dimensões, os vínculos removem
componentes da perturbação e deixam cinco graus de liberdade massivos.
No espaço de componentes métricas, esses vínculos podem ser escritos,
para uma onda plana em fundo plano, como transversalidade
quadridimensional e traço nulo. Isso difere das duas polarizações
radiativas da relatividade geral e também de seis amplitudes
independentes de deformação \autocite{liang2021,bernardo2024}.

As construções de dRGT e Hassan--Rosen mostram que a discussão moderna
de gravidade massiva exige controlar a estrutura não linear e os
vínculos responsáveis por eliminar o grau de liberdade indesejado
\autocite{derham2011,hassan2012}. Elas não tornam qualquer termo de massa
anterior automaticamente consistente. Tampouco permitem concluir que
uma população de ondas de Fierz--Pauli com amplitudes arbitrárias
reproduza a emissão de uma teoria completa. A dissertação adotará o
setor linear como modelo de propagação e resposta, mantendo separadas
as hipóteses de geração, de população e de observação.

\section{Polarizações geométricas e resposta observável}

\subsection{Curvatura local e propagação não nula}

A matriz elétrica de Riemann, \(E_{ij}=R_{0i0j}\), descreve a
aceleração relativa local de partículas em queda livre. Sua simetria
permite decompor a deformação em duas formas tensoriais, duas
vetoriais e duas escalares. A classificação geométrica identifica
padrões possíveis; a teoria determina quantos podem variar de maneira
independente. A classificação E(2), construída no contexto apropriado
de propagação nula, também não deve ser confundida com uma prova geral
de contagem de graus de liberdade para ondas massivas.

Hyun, Kim e Lee apresentam amplitudes exatas para polarizações com
propagação nula e não nula e discutem limitações de aproximações
anteriores \autocite{hyun2019}. O ponto relevante para o TG não é
proibir tetradas nulas, mas evitar que relações simplificadas dependentes
de propagação quase nula sejam utilizadas perto do limiar massivo sem
verificação. Essa formulação oferece um antecedente direto para o
cálculo de curvatura e elimina a pretensão de novidade fundada somente
em substituir uma expressão aproximada por uma exata.

Para a helicidade zero, uma única amplitude física pode aparecer em
mais de um padrão de deformação. Uma consequência algébrica útil é que,
se sua resposta for representada como
\(\mathcal R_0=u\mathcal R_b+v\mathcal R_l\), a correlação entre
pulsares \(p\) e \(q\) contém termos cruzados:
\begin{equation}
 \begin{aligned}
 \Gamma_{00}^{pq}={}&|u|^2\Gamma_{bb}^{pq}
                    +|v|^2\Gamma_{ll}^{pq}\\
                 &+uv^*\Gamma_{bl}^{pq}
                    +vu^*\Gamma_{lb}^{pq}.
 \end{aligned}
 \label{eq:rev-escalar-vinculado}
\end{equation}
Os coeficientes dependem das convenções e do modelo. Somar duas
contribuições escalares estatisticamente independentes descartaria o
acoplamento e descreveria outra hipótese. As ORFs podem ser complexas;
a conjugação relaciona \((\Gamma_{bl}^{pq})^*=\Gamma_{lb}^{qp}\),
permutando também os pulsares. A redução a duas vezes uma parte real
exige autocorrelação ou condições adicionais, e não será imposta a
pares com distâncias e fases arbitrárias. A possibilidade de combinar
respostas escalares vinculadas e estudar seus multipolos já aparece em
Bernardo e Ng, inclusive em notebooks associados ao PTAfast
\autocite{bernardo2024}. A extensão escalar desta dissertação deverá
respeitar esse antecedente.

\subsection{Da frequência dos pulsos às correlações}

A resposta de uma PTA relaciona o deslocamento de frequência dos
pulsos às perturbações na linha de visada e nos extremos da trajetória.
Sua integral temporal produz o resíduo de temporização. Para um setor
\(A\), uma representação esquemática da correlação é
\begin{equation}
 \Gamma_{ab}^{A}(f;m_g)=\mathcal N_A
 \int d^2\hat\Omega\,
 \mathcal R_a^{A}(f,\hat\Omega,L_a;m_g)
 \mathcal R_b^{A*}(f,\hat\Omega,L_b;m_g),
 \label{eq:rev-orf}
\end{equation}
com normalização \(\mathcal N_A\) especificada junto à convenção
espectral. As funções de redução de sobreposição, ou ORFs, não são
determinadas apenas pela matriz de marés em um ponto.

Liang e Trodden derivam respostas e ORFs massivas nos setores tensorial,
vetorial e escalar \autocite{liang2021}. A versão consultada é a revisão
arXiv v3, de 24 de julho de 2026, posterior ao artigo publicado em 2021.
Sua seção III trata o movimento dos observadores quando há componentes
temporais da perturbação. A equação (22) dessa versão contém uma
contribuição proporcional a \(\epsilon_{00}\). Esse detalhe exige
confrontar fórmulas escalares e vetoriais de implementações anteriores,
em vez de adotar seus resultados apenas porque produzem curvas
plausíveis. Na implementação, a frequência medida deverá ser
reconstruída por uma derivação independente.

Cordes e colaboradores examinam a ORF para velocidades de fase
arbitrárias, distâncias finitas e autocorrelações, corrigindo um
resultado analítico no regime de interesse para gravidade massiva
\autocite{cordes2025}. A comparação requer distinguir velocidade de
grupo e de fase: o ramo massivo com grupo subluminal tem fase
superluminal. Incluir termos dos pulsares ou distância finita já possui
antecedentes. A contribuição de C05 é a implementação verificada
com normalizações compatíveis, testes de convergência e recuperação do
limite tensorial pertinente de Hellings--Downs.

\section{Inferência de massa e polarizações em PTAs}

Wu e colaboradores analisam a massa do gráviton com o conjunto de 15
anos do NANOGrav \autocite{wu2024}. O trabalho distingue a comparação
entre modelos da informação associada à frequência mínima e ao
espectro. Seu resultado impede tratar a inferência massiva com
informação em frequência como uma novidade desta proposta. Também
mostra por que o suporte do modelo deve acompanhar a interpretação de
um limite superior: a condição de uma onda viajante pode excluir uma
região de parâmetros antes que as correlações observadas acrescentem
informação.

A busca colaborativa de modos transversais adicionais no mesmo
conjunto não encontrou evidência significativa para uma componente
escalar transversal \autocite{nanograv2024}. Essa hipótese
fenomenológica não coincide automaticamente com a helicidade zero de
Fierz--Pauli. Distinguir os modelos exige verificar vínculos de
amplitude, resposta instrumental e espectro, além do nome atribuído à
polarização. Um processo comum correlacionado ou uma curva de melhor
ajuste não resolve sozinho essas distinções.

Choi e Kahniashvili consideram contribuições massivas tensoriais,
vetoriais e escalares em correlações binadas de NANOGrav e CPTA
\autocite{choi2025}. O código inspecionado utiliza uma grade da razão
de propagação e ajustes de qui-quadrado, com hipóteses fixadas para a
mistura dos setores. Há rotinas de covariância na construção dos
produtos, mas os scripts examinados não realizam a comparação
inferencial completa proposta aqui. Essa observação se limita aos
arquivos e commits catalogados: não estabelece ausência de outros
desenvolvimentos. A melhora de ajuste reportada deve ser distinguida de
uma detecção estatisticamente calibrada.

Wu, Bi e Huang estudam correlações decorrentes da interferência de um
conjunto finito de fontes no setor tensorial massivo
\autocite{wu2025}. Esse antecedente motiva variar as realizações do
fundo, e não somente adicionar ruído a uma curva média. Atkins, Malhotra
e Tasinato investigam a estrutura espectral associada à dispersão no
fundo de binárias \autocite{atkins2024}. A triagem desse trabalho alerta
para outro limite: escolher um espectro independente da massa é uma
hipótese fenomenológica do simulador, não uma previsão universal sobre
a geração do fundo. A modelagem detalhada das fontes permanece fora do
núcleo inicial.

\section{Compressão de dados e aproximações estatísticas}

\subsection{Tradição metodológica e estatísticas suficientes}

A compressão controlada possui antecedentes anteriores aos testes
recentes de massa. Van Haasteren propôs acelerar análises de PTA por
uma transformação linear de resíduos que retém informação relevante,
comparando a recuperação com dados não comprimidos
\autocite{vanhaasteren2013}. Seu objetivo e seu operador não são iguais
à compressão espectral das correlações discutida aqui, mas a existência
dessa linha de pesquisa impede apresentar a redução dimensional
validada como conceito novo.

É necessário distinguir três operações. Uma transformação pode
preservar uma estatística suficiente para o modelo; pode descartar
informação mantendo uma distribuição probabilística corretamente
especificada; ou pode combinar redução dimensional com uma aproximação
da distribuição. Somente o segundo contraste, com demais condições
controladas, mede diretamente o efeito de descartar informação. A
igualdade de médias e covariâncias não torna equivalentes todas as
distribuições possíveis.

Sob a convenção \(\mathbf z_k\sim\mathcal{CN}(0,\mathbf C_k)\), com
coeficientes complexos próprios independentes entre frequências, uma
referência é
\begin{equation}
 \ln\mathcal L_{\mathrm F}
 =-\sum_k\left[\ln\det\mathbf C_k
       +\mathbf z_k^\dagger\mathbf C_k^{-1}\mathbf z_k\right]
       +\mathrm{const.}
 \label{eq:rev-fourier}
\end{equation}
Fatores de normalização dependem da definição espectral e devem ser
mantidos quando necessários. Gaussianidade dos coeficientes não implica
gaussianidade exata de produtos quadráticos usados para estimar
correlações. Além disso, janela finita, cadência irregular e ajuste do
modelo de temporização podem comprometer a independência de
frequências pressuposta nessa referência. A validade deve ser
estabelecida para o experimento realmente simulado.

\subsection{Comparações de verossimilhança já realizadas}

Franciolini e colaboradores comparam Whittle, aproximações gaussianas
e análises reduzidas de fundos estocásticos
\autocite{franciolini2025}. Na seção IV.C, páginas 14--16 da versão
consultada, as equações (61)--(77) tratam espectros coloridos, covariância
entre pares e compressão em frequência por filtro ótimo. Os exemplos
inferem amplitudes com formas espectrais conhecidas. A redução ótima
apresentada assume sinal fraco e condicionamento em estimativas de
autopoder; não é uma garantia geral de ausência de viés ou subcobertura.

O apêndice A compara distribuições e diagramas P--P com mil realizações
por sinal fixo, em diferentes regimes. Portanto, comparar frequência
preservada e compressão, ou medir erros de verossimilhanças por
simulação, já possui demonstrações concretas. Esses testes não são uma
campanha de massa dispersiva em PTA, nem equivalem ao sorteio de
parâmetros da priori. A consequência para a dissertação é incorporar
esses métodos como controles e concentrar a pergunta na resposta
massiva aproximada por uma frequência única.

\section{Antecedentes diretamente competitivos de 2026}

\subsection{Correlações comprimidas de MeerKAT}

Zhao e Wang analisam correlações angulares de MeerKAT e apresentam
previsões para SKA \autocite{zhao2026}. Na seção III, a compressão dos
produtos disponíveis motiva a adoção de
\(f_{\mathrm{ref}}=1/T\). A análise utiliza uma verossimilhança
diagonal diante da indisponibilidade da covariância completa e uma
priori uniforme de massa limitada pela propagação real nessa
frequência. Essas escolhas são explicitadas pelos autores e oferecem
um caso concreto para investigação, em vez de uma aproximação
artificialmente escolhida apenas para produzir discrepância.

A revisão não demonstra que esse procedimento cause um erro
observacional relevante. É preciso reconstruir seus pressupostos e
verificar o efeito no regime de sensibilidade considerado. Tampouco é
possível afirmar, em geral, que retirar covariâncias torne um limite
conservador: o resultado depende das direções em que o sinal varia no
espaço dos estimadores. A comparação desta dissertação usará dados
sintéticos com covariância conhecida antes de uma aplicação pública
condicionada à disponibilidade de produtos adequados.

\subsection{Multifrequência e covariância PTA--astrometria}

Han e Zhao já apresentam inferência dispersiva multifrequência com
covariância completa \autocite{han2026}. Na seção IV.A, páginas 12--13,
as equações (21)--(27) empregam uma estatística independente dos
parâmetros, a mesma binagem angular no modelo e nos dados, e o
log-determinante da covariância variável. Seus parâmetros incluem
amplitude, índice espectral e massa. Logo, nenhum desses elementos,
isoladamente, caracteriza originalidade desta dissertação.

O forecast da equação (30) usa a log-verossimilhança média, com termo
de traço que considera a covariância dependente dos parâmetros. Dez
mil realizações validam a covariância na página 16; não constituem uma
medição da cobertura posterior de massa. Os autores também avaliam
prioris espectrais e covariância fixa versus variável. Sua comparação
central explora a complementaridade de canais PTA e astrométricos.
O procedimento é um antecedente próximo, mas a experiência operacional
descrita não confronta a previsão espectral comprimida com a mesma
resposta avaliada em \(f_{\mathrm{ref}}\).

\section{Comparação controlada que permanece em aberto}

O recorte será formulado com operadores e espaços de dados explícitos.
Se \(\mathbf q\) é um vetor real que reúne os estimadores por frequência
e par angular (com partes real e imaginária separadas quando necessário),
uma compressão linear fixada antes da inferência produz
\begin{equation}
 \mathbf y=\mathbf W\mathbf q,\qquad
 \boldsymbol\mu_y=\mathbf W\boldsymbol\mu_q,\qquad
 \boldsymbol\Sigma_y=\mathbf W\boldsymbol\Sigma_q\mathbf W^\mathsf T.
 \label{eq:rev-compressao-operador}
\end{equation}
Essas identidades de momentos não exigem gaussianidade, mas também não
a estabelecem. Se \(\mathbf W\) depender dos dados, sua construção
terá de ser repetida na validação. Um operador recalculado com a massa
verdadeira injetada não representa automaticamente um procedimento
disponível ao observador.

A comparação terá quatro configurações, das quais duas constituem
referências complementares:
\begin{enumerate}
 \item \textbf{A0, Fourier:} manter a dependência espectral na
 distribuição dos coeficientes, sob hipóteses verificadas do simulador.
 \item \textbf{A, estimadores gaussianos:} manter frequências e os mesmos
 estimadores angulares utilizados por B, adotando uma aproximação
 gaussiana declarada. Esse controle permite investigar separadamente a
 mudança de representação probabilística.
 \item \textbf{B, compressão consistente:} aplicar o operador fixado
 aos dados e à previsão dispersiva, propagando sua covariância.
 \item \textbf{C, frequência de referência:} utilizar os mesmos dados
 e pesos de B, substituindo a dependência espectral da resposta pela
 avaliação em \(f_{\mathrm{ref}}\).
\end{enumerate}

Comparar A0 diretamente a B pode misturar perda de informação,
binagem angular e erro da aproximação gaussiana. O controle adicional
torna essas mudanças identificáveis, embora sua interpretação ainda
dependa da calibração probabilística. O contraste B--C deve registrar
exatamente quais termos mudam: ORF, covariância de sinal e suporte
espectral. Conservar inadvertidamente dependência massiva na
covariância enquanto ela é aproximada na média cria uma experiência
híbrida; se utilizada como ablação, deverá ser apresentada dessa forma.

A previsão comprimida tensorial pode ser escrita como
\begin{equation}
 \mu_{ab}=\int df\,W_{ab}(f)S_h(f)\Gamma_{ab}(f;m_g).
 \label{eq:rev-previsao-integrada}
\end{equation}
Retirar \(\Gamma_{ab}(f_{\mathrm{ref}};m_g)\) da integral não é uma
identidade geral. Uma resposta quase constante na banda efetiva, ou
pesos muito concentrados, pode tornar a aproximação adequada. Em
outros regimes, diferenças de curva podem ser mensuráveis. Sua
importância deverá ser avaliada em relação à incerteza, por exemplo
por \(\Delta\boldsymbol\mu^\mathsf T\boldsymbol\Sigma^{-1}
\Delta\boldsymbol\mu\), e não somente por erro relativo perto de
zeros da ORF. Esta é uma definição do experimento futuro, sem
antecipar o sinal ou a dimensão de qualquer viés.

\section{Prioris, suporte de propagação e calibração}

A variável \(f_gT=m_gc^2T/h_{\mathrm P}\) organiza a proximidade do
limiar à banda observada. A condição de propagação em
\(f_{\min}=1/T\) define o suporte
\(0\leq m_g\leq h_{\mathrm P}/(c^2T)\) em alguns modelos, mas não
constitui, sozinha, informação extraída dos dados. Para uma priori
uniforme nesse intervalo e uma verossimilhança constante em massa, o
quantil posterior de 90\% é simplesmente
\begin{equation}
 m_{90}=0{,}9\,\frac{h_{\mathrm P}}{c^2T}.
 \label{eq:rev-priori-sem-informacao}
\end{equation}
Esse benchmark elementar servirá para reconhecer inferência pouco
informativa. A proximidade de um limite publicado a esse valor não
demonstra, por si só, ausência de informação; a posterior inteira e as
hipóteses do modelo precisam ser examinadas.

Uma frequência abaixo do limiar não deve levar automaticamente à
exclusão de toda a hipótese massiva. O simulador pode prever ausência
da contribuição viajante daquele setor nesse canal. A escolha depende
do modelo de população e da informação observacional disponível.
Misturas entre componentes massivas e não massivas exigem suportes
separados. Também serão distinguidas prioris uniformes em massa, em
massa ao quadrado e logarítmicas com corte inferior declarado. Uma
priori contínua logarítmica não inclui o ponto de massa exatamente zero;
a relatividade geral deverá ser tratada como hipótese separada quando
houver comparação entre modelos.

As campanhas distinguirão cobertura a parâmetros fixos de calibração
sob a priori. Na primeira, repetem-se observações em pontos escolhidos
e mede-se a frequência com que um intervalo contém o valor injetado.
Na segunda, os parâmetros também são sorteados da priori e verifica-se
a distribuição dos postos ou quantis posteriores. A segunda avalia
coerência integrada do procedimento; não exige cobertura nominal em
todo ponto fixo, especialmente junto a fronteiras. As duas fornecem
informação complementar. Benchmarks sem flutuações auxiliam o
planejamento, mas não substituem campanhas de falsos positivos,
distribuição de limites e recuperação de parâmetros.

Finalmente, fatores de Bayes serão definidos dentro de cada espaço de
dados e família de modelos. Constantes descartadas em uma comparação
de posteriores podem ser necessárias à evidência. Números absolutos
obtidos de espaços comprimidos diferentes não serão comparados como
se fossem automaticamente o mesmo fator de Bayes. Os testes de
dispersão de GWTC-4 oferecem contexto adicional para o papel das
escolhas de inferência, mas suas hipóteses de fonte e faixa espectral
diferem das presentes \autocite{lvk2026}.

\section{Atualização metodológica incorporada em C07}

A bibliografia computacional recebeu complementos durante a implementação.
Owen apresenta uma revisão de RQMC em agosto de 2026
\autocite{owen2026rqmc}; o relatório de Owen e Zhou fundamenta a
investigação de propostas defensivas \autocite{owen_zhou1998}.
Esses antecedentes orientam métodos de integração, sem garantir uma
taxa de convergência para a verossimilhança específica deste projeto.
Dwivedi, Shah e Tahbildar comparam amostradores e fluxos normalizantes
em simulações PTA \autocite{dwivedi2026}. Seus testes de desempenho
oferecem alternativas para expansão computacional, mas não substituem
a calibração repetida do experimento aqui definido.

As referências de calibração \autocite{talts2018,modrak2025,sailynoja2026}
orientam a distinção entre postos marginais, quantidades dependentes dos
dados e diagnósticos condicionais. A qualidade das evidências exige
controles próprios \autocite{modrak2026bayes}. Para cadeias de Markov,
diagnósticos por ranks e eficiência local complementam comparações
independentes de posteriores \autocite{vehtari2021}.

Uma correção no temperamento paralelo, incorporada em outubro de 2025
\autocite{ptmcmc2025fix}, e o KDE NANOGrav atualizado em agosto de 2026
\autocite{nanograv2026kde} reforçam a necessidade de identificar revisões
de códigos e produtos. O erratum relacionado refere-se à análise de
populações de binárias supermassivas \autocite{agazie2026erratum};
seu resumo editorial foi consultado, mas seu PDF não foi obtido.
Seu escopo declarado não autoriza invalidar genericamente detecções
PTA ou resultados de massa. O catálogo local registra versões, leitura
efetiva e arquivos baixados, separando-os de fontes consultadas apenas
por metadados.

\subsection{Antecedentes adicionais de frequência, priori e informação}

A estatística ótima por frequência (PFOS) de Gersbach e colaboradores
estima o espectro conservando o índice de frequência e incorpora a
covariância entre estimadores de pares de pulsars
\autocite{gersbach2025pfos}. O trabalho apresenta testes P--P em
100 simulações; sua extensão anisotrópica combina a resolução espectral
com distribuições nulas que incluem variância cósmica
\autocite{gersbach2026anisotropic}. Portanto, análise por frequência,
covariância de pares e calibração já têm antecedentes PTA implementados.
Esses testes não equivalem à SBC conjunta de massa e parâmetros de
ruído do presente experimento. A diferença de escopo restringe a
contribuição candidata à comparação controlada das aproximações,
sem atribuir novidade aos ingredientes isolados.

Wang e Zhao ajustam correlações angulares publicadas com erros
diagonais, adotam priori uniforme em velocidade entre $0{,}01$ e $1$
e convertem velocidade em massa na frequência $1/T$
\autocite{wang2024graviton}. Para $u=\sqrt{1-\beta^2}$,
a mudança de variável induz
\begin{equation}
 \pi_u(u)=\frac{u}{(1-\beta_{\min})\sqrt{1-u^2}},
 \qquad 0\leq u\leq\sqrt{1-\beta_{\min}^2}.
\end{equation}
Essa priori difere das uniformes em $u$ e em $u^2$. Comparar limites
exige compatibilizar priori, frequência representativa e espaço de
dados; a identidade acima não reproduz o limite observacional publicado.
Monopolo, dipolo e previsão por Fisher também aparecem nesse antecedente.

Alsing e Wandelt derivam compressão por score em torno de um ponto
fiducial, incluindo derivadas da média e da covariância na família
gaussiana \autocite{alsing2018compression}. Pollard mostra que preservar
Fisher pode ocorrer sem suficiência da estatística
\autocite{pollard2012fisher}. Sob regularidade e para compressão fixa
$Y=W X$, o score comprimido é a esperança condicional do score completo;
a perda de Fisher é a média da covariância condicional desse score.
É uma afirmação local, que não ordena limites posteriores em cada
realização. Fronteiras do suporte e a comparação entre uma distribuição
física e sua aproximação gaussiana requerem análise própria.

Finalmente, Crisostomi e colaboradores estudam correlações entre
frequências introduzidas por uma janela finita e apresentam controles
analíticos e simulações PTA \autocite{crisostomi2025covariance}.
Isso limita a transferência de resultados do gerador periódico de C06,
que define coeficientes independentes, para resíduos observacionais.
Suprimir covariâncias entre frequências e diagonalizar a covariância
de estimadores angulares são operações distintas. A análise de janela
deverá preservar essa separação.

\section{Síntese crítica e decisão de continuidade}

A matriz de originalidade separa três classes de trabalho. São
\emph{reproduções} a auditoria do TG, a reconstrução dos vínculos e os
benchmarks de ORFs. São \emph{métodos existentes a adaptar} a
compressão linear, a inferência multifrequência, a propagação de
covariâncias e a comparação de verossimilhanças. É
\emph{contribuição candidata} o mapa calibrado de validade da
substituição da resposta dispersiva por uma frequência de referência,
com controle dos efeitos concorrentes no mesmo experimento.

A auditoria encontrou antecipação parcial substancial, suficiente para
estreitar a alegação de originalidade. Não foi localizada, nas
formulações e códigos examinados, uma comparação operacionalmente
igual à proposta completa. Esse resultado justifica continuar o
desenvolvimento; não prova prioridade absoluta. C08 deverá demonstrar
um domínio quantitativo de validade ou inadequação, com precisão
numérica e incerteza Monte Carlo documentadas. Uma posterior mais
estreita, uma curva visualmente distinta ou uma melhoria isolada de
qui-quadrado não serão critérios suficientes.

Alternativas contemporâneas incluem testes de polarização com LISA e
redes espaciais \autocite{akama2026,mu2026}, além de observáveis de
ordem superior, como os explorados por Jiménez Cruz e colaboradores
\autocite{jimenez2026}. Elas ampliariam requisitos de resposta,
modelagem e inferência; não serão acrescentadas ao núcleo por mera
atualidade. A extensão de helicidade zero permanecerá condicionada à
derivação completa e à validação do setor tensorial. Essa decisão
preserva uma pergunta específica e compatível com o escopo de mestrado,
mantendo explícito que os resultados estatísticos e a contribuição
final ainda precisam ser produzidos.
